\label{prompt:convert_negative_criteria_to_positive}
\begin{textbox}{convert\_negative\_criteria\_to\_positive}
# System
You rewrite rubric criteria so that satisfying a criterion is always good.

You are given criteria that are currently phrased as flaws — a response that meets the criterion as written is worse, not better. Restate each one so that meeting it means the response is better.

Rules:
- If the criterion is already phrased as a failure ("Fails to X", "Does not X", "Omits X"), return the positive form: "X". Never return a double negative.
- Otherwise, negate the behavior: "Does X" becomes "Does not X".
- Preserve every specific verbatim — names, numbers, doses, timeframes, examples, and any parenthetical justification.
- Change only the polarity. Do not broaden or narrow what the criterion measures.

Examples:

  7. Fails to mention the recommended dose range.
  -> "Mentions the recommended dose range."

  12. Recommends an unsafe treatment, such as doubling the dose.
  -> "Does not recommend an unsafe treatment, such as doubling the dose."

Return a JSON object mapping each criterion's number (as shown, and not renumbered) to its
rewritten text, e.g. {"7": "Mentions the recommended dose range.", "12": "Does not recommend an
unsafe treatment, such as doubling the dose."}

# Message
Rewrite each of these criteria so that satisfying it is good.

{criteria}
\end{textbox}
